\section{Design Pattern}
\subsection{MVVM: ModelViewViewModel}
\subsubsection{Descrizione del pattern}
Il pattern \textit{Model-View-ViewModel} separa la logica di presentazione dall'interfaccia grafica e dalla gestione dei dati.\\
La \textit{View} si occupa esclusivamente della resa grafica e dell'interazione con l'utente, il \textit{Model} incapsula dati e regole di dominio o di validazione$^G$, mentre il \textit{ViewModel} media tra i due livelli esponendo alla vista uno stato gi\`a pronto per essere utilizzato.\\
Questa suddivisione evita che la logica applicativa venga distribuita nei componenti grafici, favorendo una struttura pi\`u chiara, modulare e manutenibile.
\subsubsection{Motivazioni dell'utilizzo del pattern}
Nel progetto$^G$ SmartOrder questo pattern \`e stato adottato nel frontend$^G$ per isolare la logica di gestione delle interazioni utente dalla componente puramente visuale.\\
La scelta risulta particolarmente adatta in presenza di componenti che devono gestire stato locale, validazione$^G$ dei campi, messaggi di errore e invocazioni alle API$^G$ backend$^G$, poich\'e evita di concentrare tali responsabilit\`a direttamente nei componenti React$^G$.\\
L'adozione di MVVM migliora inoltre la testabilit\`a delle funzionalit\`a lato client, permette di riutilizzare la logica di presentazione e rende pi\`u semplice evolvere l'interfaccia senza modificare le regole sottostanti.
\subsubsection{Utilizzo del pattern nel progetto}
Nel progetto il pattern viene applicato nelle principali funzionalit\`a del frontend$^G$:
\begin{itemize}
	\item \textbf{Autenticazione}: la vista gestisce la resa dei form, mentre il \textit{ViewModel} coordina validazione$^G$, invio dei dati e aggiornamento dello stato.
	\item \textbf{Chat}: la componente grafica mostra i messaggi e raccoglie l'input dell'utente, mentre il \textit{ViewModel} governa il flusso conversazionale e le chiamate ai servizi.
	\item \textbf{Storico ordini}: la vista presenta l'elenco e i dettagli degli ordini, mentre il \textit{ViewModel} gestisce caricamento, trasformazione dei dati e stato di errore o attesa.
\end{itemize}
In questo modo la componente grafica rimane semplice e dichiarativa, mentre la logica di presentazione resta centralizzata in moduli dedicati e separata dalle regole applicative.

\subsection{Dependency Injection}
\subsubsection{Descrizione del pattern}
La \textit{Dependency Injection} \`e un pattern che prevede che le dipendenze necessarie a una classe o a un modulo non vengano create internamente, ma fornite dall'esterno al momento dell'inizializzazione o dell'invocazione.\\
In questo modo un componente dipende da astrazioni o contratti ben definiti, invece che da implementazioni concrete istanziate direttamente al proprio interno.\\
Il risultato \`e un accoppiamento minore tra le parti del sistema e una maggiore facilit\`a nella sostituzione, configurazione e verifica$^G$ delle dipendenze.
\subsubsection{Motivazioni dell'utilizzo del pattern}
Nel progetto$^G$ SmartOrder la \textit{dependency injection} \`e stata adottata per mantenere coerente l'architettura esagonale$^G$ e ridurre la dipendenza diretta tra servizi applicativi, adapter e repository.\\
Questa scelta permette di costruire catene di componenti in cui ogni elemento riceve solo le collaborazioni di cui ha bisogno, rendendo pi\`u chiaro il flusso delle responsabilit\`a.\\
L'approccio risulta inoltre vantaggioso per il testing, poich\'e consente di sostituire facilmente implementazioni reali con mock o stub, e semplifica l'evoluzione del sistema quando una tecnologia esterna cambia.
\subsubsection{Utilizzo del pattern nel progetto}
L'utilizzo del pattern è evidente soprattutto nel backend$^G$, dove i controller FastAPI ricevono i servizi applicativi tramite funzioni dedicate di costruzione.\\
Ad esempio, nei moduli \textit{RecordingController} e \textit{StoricoController} le funzioni \textit{get\_service} istanziano e collegano repository, adapter e service, per poi fornire al router un oggetto gi\`a configurato.\\
Anche i service applicativi, come \textit{RecordingService}, \textit{HistoryService}, \textit{ChatService} e i servizi di gestione del carrello, ricevono tramite costruttore le rispettive porte o dipendenze infrastrutturali, senza crearle autonomamente.\\
Questo consente di preservare l'indipendenza della logica applicativa rispetto ai dettagli tecnologici e di mantenere la composizione del sistema centralizzata e controllata.

\subsection{Adapter}
\subsubsection{Descrizione del pattern}
Il pattern \textit{Adapter} consente di mettere in comunicazione due componenti con interfacce incompatibili, introducendo un elemento intermedio che traduce richieste e dati da un formato all'altro.\\
In un'architettura a porte e adattatori, l'adapter implementa una porta definita dal dominio e la collega a una specifica tecnologia esterna, come un database$^G$, un servizio remoto o una libreria di terze parti.\\
Questo permette al nucleo applicativo di dipendere solo da contratti astratti, delegando agli adattatori i dettagli di integrazione.
\subsubsection{Motivazioni dell'utilizzo del pattern}
SmartOrder fa ampio uso del pattern Adapter perché deve interagire con più componenti infrastrutturali, tra cui il database$^G$, servizi di trascrizione audio, sistemi di validazione$^G$ e componenti di supporto alla conversazione.\\
L'introduzione degli adapter consente di evitare che il dominio o i service applicativi conoscano direttamente i dettagli implementativi di tali strumenti.\\
Questa scelta migliora la manutenibilità del sistema, favorisce la sostituzione di una tecnologia con un'altra e rende più chiara la distinzione tra logica di business e logica di integrazione.
\subsubsection{Utilizzo del pattern nel progetto}
Il pattern è utilizzato in modo sistematico nel backend$^G$ in coerenza con l'architettura esagonale$^G$.\\
Ad esempio, gli adapter di persistenza implementano le porte verso il database$^G$, mentre altri adapter collegano il sistema ai servizi impiegati per trascrizione, conversazione e validazione$^G$.\\
In tutti i casi l'adapter riceve richieste formulate secondo il linguaggio del dominio applicativo e le traduce in chiamate concrete verso repository o servizi esterni.\\
Il risultato \`e una struttura pi\`u modulare, in cui le componenti centrali del sistema restano indipendenti dalle scelte tecnologiche adottate.